%% toc: P3 41 41.2
\section{Nilpotent groups of class two and further group examples}
\label{sec:p03-w03-p03-074-nilpotent}

Nilpotent groups of class two are the first groups beyond the abelian ones,
and they already carry the full logical strength of the commutative rings
from which they can be built. This section gives the first-order
description of nilpotency of a fixed class, the computing rules for class
two, and the Heisenberg group \(\mathrm H(R)\) of a commutative ring
\(R\). The main result shows that \(R\) and \(\mathrm H(R)\) have
many-one equivalent complete theories, uniformly in \(R\): the ring is
recovered inside the group by commutators and centralizers from two
parameters, and the parameters are then removed by a first-order
condition that recognizes the good pairs. The section closes with a group
example of the opposite kind, a family of finite groups showing that
nilpotency of unbounded class is not first-order.

The inputs are a group, or a commutative ring with identity, and a
first-order sentence. The conclusions are: a single universal sentence
axiomatizing each fixed nilpotency class; a computable translation of ring
sentences into group sentences and back; the undecidability of the theory
of \(\mathrm H(\mathbb Z)\) and the decidability of the theory of
\(\mathrm H(K)\) whenever the theory of the field \(K\) is decidable; and
the failure of every first-order axiomatization of the class of all
nilpotent groups. Interpretations and the truth lemma are those of
\cref{def:p03-imaginaries-interpretation} and
\cref{thm:p03-imaginaries-quotient-truth}; the transfer of decidability
along an interpretation is \cref{thm:p03-w03-p03-022-transfer};
ultraproducts and \L o\'s's theorem are \cref{sec:types-ultraproducts-05} and
\cref{thm:los}. Modules and abelian groups, together with the accounts of
the appendix of examples in Hodges's \emph{Model Theory} and of Eklof's
chapter on applications to algebra, belong to the preceding section; nothing
below depends on them. Every result below is proved in full in the annex.

\subsection{Commutators, centres and the class laws}
\label{sec:p03-w03-p03-074-class-laws}

\begin{hbconvention}{def:p03-w03-p03-074-conventions}{Group signature and commutators}
Groups are structures for the signature \(\{\cdot,{}^{-1},1\}\) and
satisfy the usual three group axioms. For elements or terms \(x,y\) put
\[
  x^y=y^{-1}xy,\qquad [x,y]=x^{-1}y^{-1}xy .
\]
Left-normed commutators are defined by \([x_1]=x_1\) and
\([x_1,\ldots,x_{k+1}]=[[x_1,\ldots,x_k],x_{k+1}]\) for \(k\ge1\). For
\(k\ge1\) the \emph{class law of weight \(k\)} is the universal sentence
\[
  \lambda_k:\qquad \forall x_1\cdots\forall x_k\;[x_1,\ldots,x_k]=1 .
\]
The \emph{centre} of a group \(G\) is
\(Z(G)=\{g\in G:gh=hg\text{ for all }h\in G\}\), and the
\emph{centralizer} of \(a\in G\) is \(C_G(a)=\{g\in G:ga=ag\}\). In
any group \(gh=hg\) holds exactly when \([g,h]=1\).
\index{commutator!left-normed}\index{centre of a group}
\end{hbconvention}

\begin{hbdefinition}{def:p03-w03-p03-074-upper-central}{Upper central series and nilpotency class}
For a group \(G\) define subsets \(Z_i(G)\) by \(Z_0(G)=\{1\}\) and
\[
  Z_{i+1}(G)=\{g\in G:[g,h]\in Z_i(G)\text{ for every }h\in G\}.
\]
Thus \(Z_1(G)=Z(G)\). The group \(G\) is \emph{nilpotent of class at most
\(c\)} (for a natural number \(c\)) when \(Z_c(G)=G\); it is
\emph{nilpotent} when it is nilpotent of class at most some \(c\), and its
\emph{class} is then the least such \(c\). The trivial group is the only
group of class \(0\); the nontrivial abelian groups are the groups of
class \(1\).
\index{nilpotent group}\index{upper central series}
\end{hbdefinition}

\begin{hbtheorem}{thm:p03-w03-p03-074-class-law}{proved=annex:p03-w03-p03-074-class-law}{Each nilpotency class is axiomatized by one universal law}
Let \(G\) be a group and \(c\ge0\).
\begin{enumerate}
  \item \(Z_c(G)=G\) if and only if \(G\models\lambda_{c+1}\).
  \item If \(G\models\lambda_{c+1}\) then \(G\models\lambda_{c+2}\);
        hence a group of class at most \(c\) has class at most every
        \(c'\ge c\).
  \item The groups of class at most \(c\) are exactly the models of the
        group axioms together with the single universal sentence
        \(\lambda_{c+1}\). This class is closed under subgroups,
        homomorphic images, arbitrary direct products and ultraproducts.
\end{enumerate}
\end{hbtheorem}

The law form makes the first-order character of a fixed class explicit;
Burris and Sankappanavar set the expressibility of nilpotency of class
\(k\) as an exercise \cite[Chapter~V, \S1, Exercise~5(c), p.~201]{BurrisSankappanavar2012}.
For class two the theorem says that \(G\) has class at most \(2\) exactly
when every commutator lies in the centre, the law being
\(\forall x\forall y\forall z\,[[x,y],z]=1\).

\begin{hblemma}{thm:p03-w03-p03-074-centre-step}{proved=annex:p03-w03-p03-074-centre-step}{Passing to the central quotient lowers the weight by one}
Let \(G\) be a group and \(k\ge1\).
\begin{enumerate}
  \item \(G\models\lambda_{k+1}\) if and only if
        \([g_1,\ldots,g_k]\in Z(G)\) for all \(g_1,\ldots,g_k\in G\).
  \item If \(\rho:G\to H\) is a surjective homomorphism of groups whose
        kernel is \(Z(G)\), then \(G\models\lambda_{k+1}\) if and only if
        \(H\models\lambda_k\).
\end{enumerate}
\end{hblemma}

\begin{hblemma}{thm:p03-w03-p03-074-class-two-rules}{proved=annex:p03-w03-p03-074-class-two-rules}{Computing rules in class two}
In every group, \([x,y]^{-1}=[y,x]\), \(yx=xy[y,x]\) and
\([xy,z]=[x,z]^y[y,z]\). If \(G\models\lambda_3\), then for all
\(x,y,z\in G\), every integer \(m\) and every natural number \(n\):
\begin{enumerate}
  \item \([x,y]\in Z(G)\);
  \item \([xy,z]=[x,z][y,z]\) and \([x,yz]=[x,y][x,z]\);
  \item \([x^m,y]=[x,y]^m=[x,y^m]\);
  \item \((xy)^n=x^ny^n[y,x]^{n(n-1)/2}\).
\end{enumerate}
\end{hblemma}

Items 2 and 3 say that in class two the commutator is a biadditive map
from \(G\times G\) into the abelian group \(Z(G)\), alternating in the
sense that \([x,x]=1\). The examples below make this map concrete.

\subsection{Heisenberg groups over commutative rings}
\label{sec:p03-w03-p03-074-heisenberg}

\begin{hbdefinition}{def:p03-w03-p03-074-heisenberg}{The Heisenberg group of a ring}
Let \(R\) be a commutative ring with identity; \(R\) may be the zero
ring. The \emph{Heisenberg group} \(\mathrm H(R)\) has domain \(R^3\) and
product
\[
  (a,b,c)(a',b',c')=(a+a',\,b+b',\,c+c'+ab').
\]
Write \(\pi(a,b,c)=(a,b)\in R^2\), \(\zeta(c)=(0,0,c)\), and for
\(g=(a,b,c)\) and \(h=(a',b',c')\)
\[
  \Delta(g,h)=ab'-a'b,
\]
the determinant of the pair \(\pi g,\pi h\). Put
\(e_1=(1,0,0)\) and \(e_2=(0,1,0)\). The triple \((a,b,c)\) is the
upper unitriangular matrix with entries \(a,c\) in the first row and
\(b\) in the second, and the product is matrix multiplication.
\index{Heisenberg group}
\end{hbdefinition}

\begin{hblemma}{thm:p03-w03-p03-074-heisenberg-basic}{proved=annex:p03-w03-p03-074-heisenberg-basic}{Structure of the Heisenberg group}
Let \(R\) be a commutative ring with identity.
\begin{enumerate}
  \item \(\mathrm H(R)\) is a group with identity \((0,0,0)\) and
        \((a,b,c)^{-1}=(-a,-b,-c+ab)\); \(\pi\) is a surjective
        homomorphism onto the additive group \(R^2\), and \(\zeta\) is an
        injective homomorphism from the additive group of \(R\).
  \item \([g,h]=\zeta(\Delta(g,h))\) for all \(g,h\in\mathrm H(R)\).
  \item \(Z(\mathrm H(R))=\zeta(R)=\ker\pi\), and
        \(C_{\mathrm H(R)}(g)=\{h:\Delta(g,h)=0\}\).
  \item \(\mathrm H(R)\models\lambda_3\). If \(R\) is not the zero
        ring then \(\mathrm H(R)\) is not abelian, so its class is
        exactly \(2\).
  \item \((a,b,c)^n=\bigl(na,\,nb,\,nc+\tfrac{n(n-1)}2ab\bigr)\) for every
        natural number \(n\), where \(\tfrac{n(n-1)}2ab\) is that natural
        number times \(ab\).
\end{enumerate}
\end{hblemma}

The ring is recovered from the group on the centre. Addition is the group
product there; multiplication needs two parameters, which fix a
coordinate system.

\begin{hbdefinition}{def:p03-w03-p03-074-ring-formulas}{Defining the ring on the centre}
Let \(\Lang_{\mathrm{rg}}=\{0,1,+,\cdot\}\). In the group signature,
with free variables \(x,y\) for the parameters, put
\[\begin{aligned}
  \mathrm{Ctr}(z)&:=\forall w\,(zw=wz),\\
  \mathrm{Mul}(x,y;z_1,z_2,z_3)&:=\exists s\,\exists t\,
     \bigl([s,x]=1\land[t,y]=1\land[s,y]=z_1\\
     &\qquad\land[x,t]=z_2\land[s,t]=z_3\bigr),\\
  \beta(x,y)&:=\forall z\,\bigl(\mathrm{Ctr}(z)\to
       \exists s\,([s,x]=1\land[s,y]=z)\bigr)\\
     &\qquad\land\forall z\,\bigl(\mathrm{Ctr}(z)\to
       \exists t\,([t,y]=1\land[x,t]=z)\bigr).
\end{aligned}\]
Let \(c_x,c_y\) be new constants. The data \(I\) assign to the one sort
of \(\Lang_{\mathrm{rg}}\) the domain \(\mathrm{Ctr}(z)\) and literal
equality, to \(0\) the graph \(z=1\), to \(1\) the graph
\(z=[c_x,c_y]\), to \(+\) the graph \(z_3=z_1z_2\), and to \(\cdot\)
the graph \(\mathrm{Mul}(c_x,c_y;z_1,z_2,z_3)\). For an
\(\Lang_{\mathrm{rg}}\)-sentence \(\sigma\), let \(\sigma^I(x,y)\) be
the translation of \cref{def:p03-imaginaries-translation} with \(c_x,c_y\)
replaced by fresh variables \(x,y\), and put
\[
  \sigma^{\sharp}:=\exists x\,\exists y\,\bigl(\beta(x,y)\land\sigma^I(x,y)\bigr).
\]
\end{hbdefinition}

\begin{hbtheorem}{thm:p03-w03-p03-074-ring-recovery}{proved=annex:p03-w03-p03-074-ring-recovery}{The ring inside its Heisenberg group}
Let \(R\) be a commutative ring with identity and \(x,y\in\mathrm H(R)\)
with \(\delta=\Delta(x,y)\) a unit of \(R\). Then \(I\) is an
interpretation of \(\Lang_{\mathrm{rg}}\)-structures in
\((\mathrm H(R),x,y)\), with \(c_x,c_y\) naming \(x,y\) and source theory
\(\operatorname{Th}(\mathrm H(R),x,y)\). In \(I(\mathrm H(R),x,y)\) the
product of \(\zeta(r_1)\) and \(\zeta(r_2)\) is
\(\zeta(r_1r_2\delta^{-1})\), the identity is \(\zeta(\delta)\), and
\[
  \theta_{x,y}:\zeta(r)\longmapsto r\delta^{-1}
\]
is an isomorphism of \(I(\mathrm H(R),x,y)\) onto \(R\). Consequently,
for every \(\Lang_{\mathrm{rg}}\)-sentence \(\sigma\),
\(R\models\sigma\) if and only if \(\mathrm H(R)\models\sigma^I(x,y)\).
\end{hbtheorem}

\begin{hblemma}{thm:p03-w03-p03-074-unit-criterion}{proved=annex:p03-w03-p03-074-unit-criterion}{The good parameter pairs are definable}
Let \(R\) be a commutative ring with identity and
\(x,y\in\mathrm H(R)\). Then \(\mathrm H(R)\models\beta(x,y)\) if and
only if \(\Delta(x,y)\) is a unit of \(R\). In particular
\(\mathrm H(R)\models\beta(e_1,e_2)\).
\end{hblemma}

\begin{hbtheorem}{thm:p03-w03-p03-074-equivalence}{proved=annex:p03-w03-p03-074-equivalence}{A commutative ring and its Heisenberg group have equivalent theories}
Let \(R\) be a commutative ring with identity.
\begin{enumerate}
  \item For every \(\Lang_{\mathrm{rg}}\)-sentence \(\sigma\),
        \(R\models\sigma\) if and only if
        \(\mathrm H(R)\models\sigma^{\sharp}\), and \(\sigma\mapsto\sigma^{\sharp}\)
        is computable and independent of \(R\).
  \item \(\mathrm H(R)\) is interpreted in \(R\) without parameters, in
        dimension \(3\), by the domain formula \(a=a\), coordinatewise
        equality, the identity graph
        \(z_1=0\land z_2=0\land z_3=0\), the inverse graph
        \(a+z_1=0\land b+z_2=0\land c+z_3=a\cdot b\), and the product
        graph \(z_1=a+a'\land z_2=b+b'\land z_3=(c+c')+a\cdot b'\);
        the identity map of \(R^3\) is the isomorphism.
  \item Hence \(\operatorname{Th}(R)\) and \(\operatorname{Th}(\mathrm H(R))\)
        are each many-one reducible to the other by computable maps that do
        not depend on \(R\). One is decidable if and only if the other is.
\end{enumerate}
\end{hbtheorem}

\begin{hbcorollary}{thm:p03-w03-p03-074-examples}{proved=annex:p03-w03-p03-074-examples}{Decidable and undecidable Heisenberg groups}
\begin{enumerate}
  \item The theory of \(\mathrm H(\mathbb Z)\) is undecidable.
  \item If \(K\) is a field whose theory in \(\Lang_{\mathrm{rg}}\) is
        decidable, the theory of \(\mathrm H(K)\) is decidable; this
        applies to every finite field.
\end{enumerate}
\end{hbcorollary}

Item 1 uses the undecidability of the ring of integers,
\cref{thm:p02-w02c-p02-003-integers}. For item 2, the decidability of the
theories of algebraically closed fields of fixed characteristic and of real
closed ordered fields is the subject of
\cref{sec:p03-w03-p03-022-procedures}; whenever such a theory is available
as a decidable complete theory of a field \(K\), item 2 transfers it to
\(\mathrm H(K)\).

\begin{hbexample}{def:p03-w03-p03-074-worked}{Multiplying integers with commutators}
In \(\mathrm H(\mathbb Z)\) take the parameters \(e_1,e_2\), so
\(\delta=1\cdot1-0\cdot0=1\). To compute \(2\cdot3\), the centralizer
condition \([s,e_1]=1\) forces \(\pi s=(\alpha,0)\), and \([s,e_2]=\zeta(2)\)
forces \(\alpha=2\); take \(s=(2,0,0)\). Similarly \(t=(0,3,0)\) satisfies
\([t,e_2]=1\) and \([e_1,t]=\zeta(3)\). Then
\([s,t]=\zeta(2\cdot3-0\cdot0)=\zeta(6)\), so
\(\mathrm{Mul}(e_1,e_2;\zeta(2),\zeta(3),\zeta(6))\) holds, and
\(\theta_{e_1,e_2}\) sends these to \(2,3,6\).

Now take \(x=e_2\) and \(y=e_1\), so \(\delta=0\cdot0-1\cdot1=-1\). The
witness \(s\) must satisfy \(\pi s=(0,\alpha)\) and
\([s,y]=\zeta(-\alpha)\), and \(t\) must satisfy \(\pi t=(\beta,0)\) and
\([x,t]=\zeta(-\beta)\). For \(z_1=\zeta(2)\) and \(z_2=\zeta(3)\) this
gives \(\alpha=-2\), \(\beta=-3\) and
\([s,t]=\zeta(0\cdot0-(-2)(-3))=\zeta(-6)\). The new ring has identity
\(\zeta(-1)\) and product \(\zeta(2)\ast\zeta(3)=\zeta(-6)\), in
agreement with \(r_1r_2\delta^{-1}=6\cdot(-1)\); the isomorphism
\(\theta_{x,y}\) is \(\zeta(r)\mapsto-r\) and sends the three elements to
\(-2,-3,6\).

For powers, \cref{thm:p03-w03-p03-074-heisenberg-basic} gives
\(g^p=(0,0,\tfrac{p(p-1)}2ab)\) in \(\mathrm H(\mathbb F_p)\) for
\(g=(a,b,c)\). For an odd prime \(p\) the integer \(\tfrac{p(p-1)}2\) is
a multiple of \(p\), so every element has order dividing \(p\). For
\(p=2\) the element \((1,1,0)\) has square \((0,0,1)\neq1\), so
\(\mathrm H(\mathbb F_2)\) has an element of order \(4\).
\end{hbexample}

\begin{hbwarning}{def:p03-w03-p03-074-bad-pairs}{Pairs with a non-unit determinant}
The unit hypothesis cannot be dropped. In \(\mathrm H(\mathbb Z)\) take
\(x=(2,0,0)\) and \(y=e_2\), so \(\Delta(x,y)=2\). An element \(t\) with
\([t,y]=1\) has \(\pi t=(0,\beta)\), and then \([x,t]=\zeta(2\beta)\).
The relation \(\mathrm{Mul}(x,y;z_1,z_2,z_3)\) therefore holds for no
\(z_3\) when \(z_2=\zeta(1)\), so the graph of multiplication is not
total and \(I\) is not an interpretation in \((\mathrm H(\mathbb Z),x,y)\).
The formula \(\beta(x,y)\) fails at the same witness \(z=\zeta(1)\), which
is why \(\sigma^{\sharp}\) quantifies only over pairs satisfying
\(\beta\). The construction also uses commutativity of \(R\): the unit
criterion and the formula for the product rest on determinant identities
that hold in commutative rings.
\end{hbwarning}

\subsection{Unbounded class is not first-order}
\label{sec:p03-w03-p03-074-unbounded}

\begin{hbdefinition}{def:p03-w03-p03-074-dihedral}{Dihedral groups as semidirect products}
For an integer \(m\ge1\) let \(G_m\) have domain
\(\mathbb Z/m\mathbb Z\times\{1,-1\}\) and product
\[
  (a,\varepsilon)(b,\eta)=(a+\varepsilon b,\ \varepsilon\eta).
\]
It is a group with identity \((0,1)\) and
\((a,\varepsilon)^{-1}=(-\varepsilon a,\varepsilon)\); for \(m\ge3\) it is
the dihedral group of order \(2m\), with rotation \((1,1)\) and
reflection \((0,-1)\).
\index{dihedral group}
\end{hbdefinition}

\begin{hblemma}{thm:p03-w03-p03-074-dihedral-class}{proved=annex:p03-w03-p03-074-dihedral-class}{Dihedral two-groups have every class}
For every \(k\ge1\) the group \(G_{2^k}\) is nilpotent of class exactly
\(k\). For \(k\ge2\), \(Z(G_{2^k})=\{(0,1),(2^{k-1},1)\}\) and reduction
modulo \(2^{k-1}\) in the first coordinate is a surjective homomorphism
\(G_{2^k}\to G_{2^{k-1}}\) with kernel \(Z(G_{2^k})\).
\end{hblemma}

\begin{hbtheorem}{thm:p03-w03-p03-074-not-elementary}{proved=annex:p03-w03-p03-074-not-elementary}{Nilpotency is not first-order}
Let \(\mathcal U\) be a nonprincipal ultrafilter on the set of positive
integers. The ultraproduct \(\prod_{k\ge1}G_{2^k}/\mathcal U\) is a group
that is not nilpotent, although every factor is nilpotent. Consequently no
set of first-order sentences in the group signature has as its models
exactly the nilpotent groups.
\end{hbtheorem}

\begin{hbwarning}{def:p03-w03-p03-074-fixed-class-only}{Fixed class versus unbounded class}
\Cref{thm:p03-w03-p03-074-class-law} and
\cref{thm:p03-w03-p03-074-not-elementary} do not conflict. Each class
bound is one universal law, and the nilpotent groups form the union of the
increasing sequence of these elementary classes. The union is not closed
under ultraproducts, because the ultraproduct satisfies the negation of each
law \(\lambda_{c+1}\) as soon as almost all factors do. The existence of
a nonprincipal ultrafilter on an infinite set is
\cref{lem:nonprincipal-ultrafilter}, which assumes the Ultrafilter
Lemma.
\end{hbwarning}

\paragraph{Relationships and sources.}
The interpretation of \cref{thm:p03-w03-p03-074-ring-recovery} is a
quotient-free instance of \cref{def:p03-imaginaries-interpretation} with
named parameters, as permitted by \cref{def:p03-imaginaries-standing}; the
parameter-free translation \(\sigma^{\sharp}\) replaces the parameters by
a first-order description of all admissible pairs. The decidability
consequences are instances of \cref{thm:p03-w03-p03-022-transfer}, and
the failure of axiomatizability uses only \cref{thm:los}. The
expressibility of a fixed nilpotency class is posed as an exercise in
\cite[Chapter~V, \S1, Exercise~5(c), p.~201]{BurrisSankappanavar2012}; the
other statements and proofs here are derived directly in the annex.
